\section*{ЗАКЛЮЧЕНИЕ}
\addcontentsline{toc}{section}{ЗАКЛЮЧЕНИЕ}

Преимущества использования веб-платформ для поиска работы заключаются в их универсальности, масштабируемости и возможности объединения всех участников рынка труда в единой экосистеме. Платформа обеспечивает эффективное взаимодействие между соискателями, работодателями и HR-специалистами, предоставляя удобные инструменты для поиска, отбора и трудоустройства. 
  
Современный рынок труда активно использует информационные технологии для повышения эффективности процессов найма, создавая специализированные платформы, которые автоматизируют подбор кандидатов, предоставляют расширенные возможности поиска и обеспечивают удобное управление вакансиями и резюме. Для реализации этих задач было разработано веб-приложение на основе современных технологий с использованием реляционной базы данных, обеспечивающее полный цикл взаимодействия участников рынка труда.

Основные результаты работы:

\begin{enumerate}
\item Проведён анализ предметной области. Выявлена необходимость разработки комплексной платформы поиска работы с поддержкой создания резюме, публикации вакансий, управления компаниями и системы откликов.
\item Разработана концептуальная модель веб-приложения. Разработана модель данных системы. Определены требования к функциональности и пользовательскому интерфейсу.
\item Осуществлено проектирование веб-приложения. Разработана архитектура системы с разделением на клиентскую и серверную части. Спроектирован пользовательский интерфейс, обеспечивающий удобство использования для различных категорий пользователей.
\item Реализовано и протестировано веб-приложение. Проведено модульное и системное тестирование.
\end{enumerate}

Все требования, объявленные в техническом задании, были полностью реализованы, все задачи, поставленные в начале разработки проекта, были также решены.

Готовый рабочий проект представляет собой полнофункциональную веб-платформу поиска работы, обеспечивающую эффективное взаимодействие между соискателями и работодателями. Приложение готово для коммерческого использования и может быть развёрнуто на сервере для публичного доступа с возможностью дальнейшего масштабирования и добавления новых функций.
